
\chapter[APÊNDICE A -- PLANO DE TRABALHO: DIVISÃO DE RESPONSABILIDADES]{{\centering APÊNDICE A \textendash{} PLANO DE TRABALHO: DIVISÃO DE RESPONSABILIDADES}}  \label{ap:plano}

Este apêndice detalha a divisão de responsabilidades entre os dois autores, conforme a pré-proposta aprovada e o plano de trabalho entregue à Coordenação do Curso de Engenharia de Computação do Centro Federal de Educação Tecnológica de Minas Gerais, Unidade Leopoldina.

\section{Atividades Desenvolvidas em Regime Colaborativo}

As atividades listadas a seguir foram conduzidas de forma conjunta por Igor Lamoia Queiroz e Victor de Souza Vilela da Silva, com participação efetiva de ambos os integrantes nas etapas de especificação, implementação, integração e revisão:

\begin{itemize}
    \item Revisão bibliográfica sobre interpretadores, compiladores e ferramentas educacionais relacionadas;
    \item Definição dos requisitos funcionais e não funcionais e da arquitetura geral do sistema;
    \item Participação na implementação e validação do núcleo do compilador, incluindo análise léxica, análise sintática, geração e execução de código intermediário;
    \item Especificação dos mecanismos de injeção dinâmica de regras gramaticais e integração desses mecanismos à experiência de personalização da IDE;
    \item Integração do núcleo do interpretador com a plataforma web e definição dos contratos de comunicação entre os módulos;
    \item Revisão mútua de código e tomada de decisões arquiteturais nas etapas de refatoração e evolução do sistema.
\end{itemize}

\section{Eixos de Aprofundamento Individual}

A partir da base técnica construída em conjunto, cada autor concentrou seu aprofundamento acadêmico em eixos específicos, conforme estabelecido na pré-proposta.

\subsection{Igor Lamoia Queiroz (este documento)}

\begin{itemize}
    \item Detalhamento da arquitetura da aplicação web (Next.js, Pages Router) e da integração com o Monaco Editor para edição de código com realce de sintaxe configurável;
    \item Documentação da customização dinâmica de \textit{tokens} e sintaxe por meio do assistente de personalização da linguagem (\textit{wizard}), incluindo o modelo de etapas, a lógica de validação preventiva e o mecanismo de pré-visualização em tempo real;
    \item Análise da experiência de uso da plataforma, dos contextos React que governam o estado da aplicação e das mediações pedagógicas oferecidas pela interface ao estudante iniciante;
    \item Descrição da execução das análises léxica e sintática com o código produzido na IDE e da apresentação do código intermediário no terminal integrado;
    \item Estudos de caso centrados na usabilidade, no engajamento dos estudantes e no \textit{feedback} coletado em ambiente educacional.
\end{itemize}

\subsection{Victor de Souza Vilela da Silva (documento complementar)}

\begin{itemize}
    \item Detalhamento das fases de análise léxica e do sistema de varredura por escâneres especializados (\textit{factory pattern});
    \item Descrição aprofundada da análise sintática por descida recursiva e da correspondência entre as produções do analisador e a gramática formal da linguagem Java-{}-;
    \item Documentação da geração de código de três endereços: instrução a instrução, tabela de temporários e tabela de rótulos;
    \item Análise do funcionamento do interpretador, da tabela de símbolos e do despacho das chamadas de sistema de entrada e saída;
    \item Desenvolvimento do \textit{back-end} da plataforma em FastAPI, incluindo a modelagem das entidades de domínio, a implementação das rotas REST, a persistência com SQLAlchemy assíncrono e o suporte a autenticação, turmas, exercícios, listas e submissões;
    \item Estudos de caso centrados na validade técnica do núcleo de tradução e no comportamento do interpretador frente a programas de teste elaborados para cobrir as construções da linguagem.
\end{itemize}
